\documentclass[12pt]{article}
\usepackage[utf8]{inputenc}
\usepackage[T1]{fontenc}
\usepackage{hyperref}
\title{Coatl\'an Zapotec Language Resources for K-12 Curriculum Development}
\author{Compiled by Claude Sonnet 5.5 (claude-sonnet-5-5) for Andr\'as Kornai}
\date{}
\begin{document}
\maketitle

\section{Introduction}
Coatl\'an Zapotec (ISO 639-3 zps; Glottocode coat1243) is an Otomanguean language of Oaxaca, Mexico, also called Western Miahuatl\'an Zapotec or San Miguel Zapotec. It has about 500 speakers. Vitality assessments conflict: Ethnologue gives EGIDS 7, while a 2013 Kornai rating says ``dead''. It must not be confused with San Vicente Coatl\'an Zapotec (zpt), which has a published New Testament. That New Testament does not belong to zps. Very little exists for zps.

\section{Primary Educational Resources}
None were found.

\section{Linguistic and Reference Materials}
\begin{itemize}
\item Language catalogue entry, which reports one primary text online and no spellchecker, keyboard or Wikipedia: \url{https://hlt.bme.hu/en/dld/language/872}
\end{itemize}

\section{Religious and Community Resources}
None were found for zps.

\section{Digital and Audio Resources}
\begin{itemize}
\item Global Recordings Network audio-visual Bible lessons (40 sections, evangelistic), with MP3 and slideshow downloads: \url{https://globalrecordings.net/en/language/zps}
\end{itemize}

\section{Curriculum Development Context}
No dictionary, grammar or corpus was found, and the language is not known to be taught in schools. A curriculum would require new documentation with the remaining speakers, who are reportedly adults only.
\end{document}